\documentclass[10 pt,usenames,dvipsnames, oneside]{article}
\usepackage{../../../modelo-ensino-medio}



\begin{document}

\begin{center}
  \begin{minipage}[l]{3cm}
\includegraphics[width=2cm]{logo}    
\end{minipage}\hfill
\begin{minipage}[r]{.8\textwidth}
 {\Large \scshape Atividade: Produção de peças}  
\end{minipage}
\end{center}
\vspace{.2cm}

\ifdefined\prof
%Habilidades da BNCC
\begin{objetivos}
\item a
\end{objetivos}

%Caixa do Para o Professor
\begin{goals}
%Objetivos específicos
\begin{enumerate}
\item Aplicar a definição de probabilidade condicional para reconhecer eventos independentes e eventos dependentes.
\end{enumerate}

\end{goals}

\bigskip
\begin{center}
{\large \scshape Atividade}
\end{center}
\fi

\begin{figure}[H]
\centering

\noindent\includegraphics[width=350bp]{maquinas}

\caption{Produção de peças}
\end{figure}



Para fins de controle de qualidade das peças produzidas em uma fábrica, foram analisadas amostras de peças produzidas por máquinas diferentes. Os resultados obtidos estão resumidos no quadro a seguir.

\begin{table}[H]
\centering
\begin{tabu} to \textwidth{|c|c|c|}
\hline
\thead
Máquina
&
Defeituosas
&
Boas
\\
\hline
$I$
&
5
&
195
\\
\hline
$II$
&
15
&
585
\\
\hline
\end{tabu}
\end{table}


Pela informação das amostras coletadas, há razão para se acreditar que a produção de defeituosas ou boas é independente do tipo de máquina utilizada na produção?

\ifdefined\prof
\begin{solucao}

Defina os eventos $A$: “peça produzida pela máquina $I”$ e $B$: “peça defeituosa”. Observe que os respectivos eventos complementares são $\overline{A}$: “peça produzida pela máquina $II$”, pois são apenas duas máquinas e $\overline{B}$: “peça boa”, pois são apenas duas classificações das peças (boas ou defituosas).

Se os eventos são independentes, deve-se ter $P(B|A)=P(B)$, por exemplo. Tem-se que $P(A)=\frac{200}{200+600}=0{,}25$, $P(B)=\frac{20}{800}=0{,}025$ e $P(B|A)=P(A\cap B)P(A)=\frac{5/800}{200/800}=\frac{1}{40}=0{,}025$. Logo, conlui-se que a produção de defeituosas é independente do tipo de máquina, pois $P(B|A)=P(B)$.

Observe que tanto faz a escolha dos eventos máquina $I$ ou máquina $II$ e peça defeituosa ou peça boa. Por exemplo,

$P(\overline{A})=\frac{3}{4}=0{,}75, P(B)=\frac{20}{800}=0{,}025$ e $P(\overline{B})=0{,}975$. Além disso, $P(\overline|B)=0{,}75=P(\overline{A}), P(A|\overline{B})=0{,}25=P(A), P(\overline{A}|\overline{B})=0{,}75=P(\overline{A})$.

De fato, essa propriedade não é um caso específico desse problema e ela é explorada na próxima atividade.

\end{solucao}
\fi

\end{document}